\chapter{Evaluation protocol}
\label{cap:protocol}

\section{Two levels and a criterion for the main question}
\label{sec:two-levels}

\todo{N1, the comparative level (do graphs from different representations
differ significantly in quality and structure?), and N2, the absolute and
validity level (does a representation induce a good graph on its own, and do
genre-based conclusions hold up against human judgements?). Say why they are
reported separately: they answer different questions, and a representation can
do well on one and badly on the other. Note that probing classifiers are a
diagnostic used inside N1 (an explanatory variable for SQ3), not a third
level.}

\subsection{Criterion for answering the main question}
\label{sec:pq-criterion}

\todo{State, verbatim and before any result is seen, the criterion that decides
the answer to PQ (``do differences between representations translate into
significant differences in graph structure and quality?''):
\begin{itemize}
  \item \textbf{Yes}, if both hold: (1) at least one comparison between
        representations of the declared family has a $\Delta$ whose 95\,\% CI
        excludes 0 after the Holm correction, in some headline quality metric
        \emph{and} some headline structure metric; and (2) the dissimilarity
        between graphs of different representations exceeds the noise floor
        measured between seeds of the same representation.
  \item \textbf{Partially}, if only one of the two holds, or if it holds for
        quality only or for structure only.
  \item \textbf{No}, otherwise.
\end{itemize}
State that the answer is given separately for the FMA test gallery and for
MTT, and that SQ5 (siamese inputs versus each other and versus classical
techniques) and SQ6 (fuzzy versus crisp) are answered with the same
statistical protocol, over their own declared comparisons.}

\section{Retrieval}
\label{sec:retrieval-protocol}

\todo{Gallery, query protocol and the exact definitions. The gallery is always
the whole corpus; \texttt{queries} only decides which queries the averages run
over, and \texttt{test} is the headline protocol once a model has been trained
on \texttt{training}. Give the formulas for R@k (raw and normalised by
$\min(1, k / |\text{relevant}|)$), MAP, nDCG@k (binary and graded) and P@k,
including the denominator of AP@k. State clearly which metrics are headline and
which are secondary: \textbf{Recall@10 (raw and normalised), MAP and nDCG@10}
are the headline quality metrics, together with agreement on the MTT triplets;
\textbf{P@10 is kept as a secondary metric}, only for continuity with the
regression figures of the archived repository. Describe the three exclusion
variants: \texttt{dedup}, the headline one, where the near-duplicate partner of
a query counts as the query itself and leaves both the ranking and the relevant
set; \texttt{raw}, with no such treatment; and \texttt{artist\_filter}, which
removes every track by the same artist.}

\section{Uncertainty}
\label{sec:uncertainty}

\todo{Percentile bootstrap with 1000 resamples. Explain the clustered variant
and why it is the one to read when comparing models: tracks by one artist are
not independent, so resampling queries understates the uncertainty. Paired
bootstrap for differences between models, so the interval is on the difference
and not on two independent means.}

\section{Agreement with human judgements}
\label{sec:triplet-protocol}

\todo{For each constraint $(s, p, n)$, a hit is $\cos(s,p) > \cos(s,n)$, with
ties counted as one half. The headline number is the unweighted mean over the
860 constraints, which is what the literature reports; the vote-weighted mean
and the subset with margin at least 2 are variants. Intervals cluster by
triplet, 337 clusters. Reference points for the table: chance is 0.5, a random
embedding is the negative control, tag Jaccard reaches about 0.57, genre reaches
0.90 on the 12\,\% of constraints where it decides at all, MFCC audio-words
reach 63.6 to 64.6\,\% \cite{Karamanolakis2016}, and metrics learned on the
constraints themselves reach 69.5 to 75.6\,\% \cite{Wolff2012}. State the
limitation plainly: 860 constraints in 337 triplets give an interval of roughly
$\pm 3$ to 5 points, which separates a baseline from a strong model but not two
close variants.}

\section{Probing}
\label{sec:probing-protocol}

\todo{Logistic and 512-unit MLP probes over frozen embeddings, the grid searched
on validation only, and the metrics per task kind. Say why touching the test
split during selection would make every published probing number optimistic.}

\section{Graph structure, hubness and robustness}
\label{sec:structure-protocol}

\todo{Community detection with several seeds and the stability of the partition
reported next to its quality; ARI, NMI, modularity, components, clustering
coefficient and silhouette. The k-occurrence distribution, its skewness,
anti-hubs, and mutual proximity. Robustness as neighbour-set Jaccard@k and
$\Delta$P@k under noise, gain, pitch shift, time stretch and cropping.}

\section{Similarity measures and classical classifiers}
\label{sec:similarity-measures-protocol}

\todo{Describe the measures compared on the classical representations and on
the best siamese model: cosine (canonical), Euclidean, Manhattan, Pearson
correlation, Mahalanobis (with Ledoit-Wolf shrinkage) and symmetric KL between
Gaussians for \texttt{gauss\_mfcc}. State that the siamese distance coincides
with the canonical cosine measure in its own space. Describe how classical
classifiers (kNN, logistic regression, SVM-RBF, random forest) are evaluated
twice: as classifiers (accuracy, macro F1 on \texttt{fma\_test}) and, through
their class-posterior vector, as graph generators alongside the learned
representations.}

\section{Fuzzy evaluation}
\label{sec:fuzzy-protocol}

\todo{Describe how the fuzzy block (F1--F4) is evaluated against its crisp
counterpart, on \texttt{mfcc}, the best siamese model and \texttt{mert}: F1
(fuzzy similarity graph) against retrieval ordered by membership degree, tie
agreement and weighted structure metrics, across a curve of $\alpha$-cuts; F2
(fuzzy inference), if confirmed by the supervisor, against cosine on
\texttt{acoustic} and against \texttt{siamese\_desc}; F3 (fuzzy genre
membership: fuzzy kNN and fuzzy c-means) against the crisp kNN and the
classical classifiers, scored with the multi-label agreement (FMA sub-genres,
MTT tags) and fuzzy ARI; F4 (model-versus-human ambiguity) as a Spearman
correlation between the human vote margin and the fuzzy membership margin on
the MTT triplets. State that every fuzzy parameter (the $\alpha$ of F1, the
rule base and membership functions of F2) is fixed on the FMA validation split
before touching test.}

\section{Invariants}
\label{sec:invariants}

The whole evaluation rests on three rules that the code enforces rather than
assumes. Rows and identifiers stay aligned, and every cross-reference goes
through the identifier, never through a position. Anything fitted, the scaler,
the PCA and the probes, is fitted on the training split alone. Every source of
randomness is seeded from the single seed the run records.
